\section{Discussion and Conclusion}\label{sec:discussion}

\subsection{What the Current Evidence Supports}

The public results support a precise and limited conclusion. The repository implements disagreement scoring and bounded debate orchestration, while the current 4,500-sample run establishes only the behavior of the local heuristic detection pipeline. It does not isolate debate and therefore provides no quantitative evidence that debate improves precision or removes false positives. The current direct evidence for debate is implementation-level: disagreement can create a structured session and preserve an auditable procedural artifact.

The system evidence is also meaningful, but should remain conservative. The validated workspace snapshot confirms that the paper is not describing an isolated algorithm detached from a deployable security stack. Authenticated gateways, signed command dispatch, audit-chain retention, and benchmark harnesses are present. However, these artifacts do not justify broad claims about production scale, full LLM latency, or externally validated rollout. The draft therefore frames \system\ as a concrete realization environment, not as proof of every historical platform claim.

\subsection{Threats to Validity}

Four threats are especially important.

First, the public corpus is synthetic and cannot estimate field detection or false-positive rates. Second, the current run uses a heuristic pipeline rather than real local and cloud models. Third, there is no external IDS/WAF baseline. Fourth, the role-specialization and debate arguments remain conceptually stronger than their empirical support because the current evidence has no controlled debate or role ablation.

The planned private labeled-alert layer addresses the first and third limitations most directly. Once that data is available, the paper should add a role-specialization ablation and a targeted error analysis on genuine alert ambiguity.

\subsection{Conclusion}

This draft reframes \system\ into a narrower, falsifiable JCST contribution. The paper does not argue that a broad security platform is itself the novelty. It proposes debate-driven multi-agent judgment as a method for distributed web threat triage and establishes that the repository contains a functioning, auditable implementation surface. Whether that method improves accuracy is deliberately left open. The next revision must add a versioned external corpus, real configured models, a controlled debate ablation, and independent error analysis before making a quantitative benefit claim.

\bmhead{Acknowledgements}

Acknowledgements are intentionally anonymized in this draft. Replace this paragraph in the non-anonymous revision with funding support, institutional assistance, infrastructure credits, or other brief acknowledgements that are appropriate for the final submission.

\section*{Declarations}

\begin{itemize}
  \item \textbf{Funding:} To be completed in the non-anonymous revision. The current draft does not attribute any grant support without repository-traceable confirmation.
  \item \textbf{Conflict of interest:} The author declares no competing interests for this anonymized draft.
  \item \textbf{Ethics approval and consent to participate:} Not applicable.
  \item \textbf{Consent for publication:} Not applicable.
  \item \textbf{Data availability:} Public evidence used in this manuscript is traceable to repository artifacts listed in Appendix~A. The private labeled-alert layer is reserved for a later revision and is not populated in the current draft.
  \item \textbf{Materials availability:} Not applicable.
  \item \textbf{Code availability:} The implementation surfaces discussed in this paper are present in the repository workspace referenced throughout the manuscript and Appendix~A.
  \item \textbf{Author contribution:} To be completed in the non-anonymous revision.
\end{itemize}
